\providecommand{\rank}{\operatorname{r}}
\providecommand{\RR}{\mathbb R}
\providecommand{\CC}{\mathbb C}
\providecommand{\KK}{\mathbb K}

% ---------- Exercise 38 ----------
\setcounter{enumi}{37}
\item {\bfseries\boldmath
Let $(S)$ be a linear system over $\RR$ in the unknowns $x,y,z,w$, with general solution
}

\[
(x,y,z,w)=(2,1,-1,0)+\alpha(1,3,0,1)+\beta(1,0,0,1),\qquad\alpha,\beta\in\RR.
\]

\begin{enumerate}[label=\alph*)]

\item Justify that $(S)$ is not homogeneous.

\begin{answer}
Write $p=(2,1,-1,0)$, $u=(1,3,0,1)$, $v=(1,0,0,1)$.
The given general solution is $p+\alpha u+\beta v$.
Every such point satisfies $z=-1$ and $x-w=2$.
Conversely, any point satisfying these two equations has the given form:
choose $\alpha=(y-1)/3$ and $\beta=w-\alpha$.

The zero vector is not a solution, since its third coordinate is not
$-1$. Every homogeneous system has the zero solution, so this system is not homogeneous.
\end{answer}

\item State, with justification, which of the following families in $\RR^4$ are general solutions of $(S)$:

\begin{enumerate}[label=\roman*)]

\item
\[
(x,y,z,w)=(3,1,-1,1)+\alpha(1,3,0,1)+\beta(1,0,0,1),\qquad\alpha,\beta\in\RR.
\]

\begin{answer}
Yes. $(3,1,-1,1)=p+v$. Replacing the particular solution by this point
amounts to replacing $\beta$ by $\beta+1$, which still runs through all real values.
\end{answer}

\item
\[
(x,y,z,w)=(1,4,-1,-1)+\alpha(1,3,0,1)+\beta(1,0,0,1),\qquad\alpha,\beta\in\RR.
\]

\begin{answer}
Yes. $(1,4,-1,-1)=p+u-2v$. Translating the two parameters by $(1,-2)$
does not change the set of solutions.
\end{answer}

\item
\[
(x,y,z,w)=(3,1,-1,0)+\alpha(1,3,0,1)+\beta(1,0,0,1),\qquad\alpha,\beta\in\RR.
\]

\begin{answer}
No. $(3,1,-1,0)$ has $x-w=3$, so the proposed family already contains
a nonsolution at $\alpha=\beta=0$.
\end{answer}
\end{enumerate}

\item Justify that the family
\[
(x,y,z,w)=(2,1,-1,0)+\alpha(3,4,-1,1)+\beta(1,0,0,1),\qquad\alpha,\beta\in\RR.
\]
is not a general solution of $(S)$.

\begin{answer}
In the proposed family the third coordinate is $-1-\alpha$.
For example, $\alpha=1$, $\beta=0$ gives $(5,5,-2,1)$, which is not a solution.
Hence that family cannot be a general solution of the given system.
\end{answer}
\end{enumerate}
